\section{\blue{Primal--Dual Formulation and Optimality Guarantees}}
\label{Section: Methodology}
% \subsection{DualHash: A Stochastic Primal-Dual Hashing Algorithm}
% Inspired by PDHG \cite{chambolleFirstOrderPrimalDualAlgorithm2011} in convex optimization and its nonconvex extension PPDG \cite{guoPreconditionedPrimalDualGradient2023}, we consider the equivalent constrained formulation of  the problem \eqref{Prop: block-wise regularized problem}:
% \begin{align}
% \min_{\bx \in \RR^{d_{\bx}}, \bB \in \RR^{n\times d}, \by \in \RR^{n\times d}} &\quad F(\bx, \bB) + h(\bY) \notag \\
% \text{s.t. } &\quad \bY = \bB \nonumber ,
% \end{align} whose corresponding dual problem is
% \begin{align}
% \label{Prop: the dual problem}
% \max_{\bLmbd \in \RR^{n\times d}} \min_{\bx \in \RR^{d_{\bx}}, \bB \in \RR^{n\times d}, \by \in \RR^{n\times d}} &\quad F(\bx, \bB) + h(\bY) + \langle \bLmbd, \bB - \bY \rangle. \tag{D}
% \end{align}
% Recalling the definition of the conjugate function \cite{Rockafellar1998VariationalA}:
% \begin{equation}
% \label{eq:fenchel_conjugate}
% h^*(\bY) := \sup_{\bx \in \RR^{n\times d}} \{ \langle \bY, \bx \rangle - h(\bx) \}, \quad \bY \in \RR^{n\times d}
% \end{equation} and observing that the minimization with respect to $\bY$ in the problem \eqref{Prop: the dual problem} can be performed independently, we then simplify the Lagrangian function in the problem \eqref{Prop: the dual problem} to
% \begin{align}
% \label{Eq: simplified_lagrangian}
% \calL(\bx, \bB, \bLmbd) := F(\bx, \bB) + \langle \bB, \bLmbd \rangle - h^*(\bLmbd).\tag{\(\calL\)}
% \end{align}

% Through this transformation, we replace $h(\bY)$ with $-h^*(\bLmbd)$ and can optimize \(h^*\) in the dual space. Since $h^*$ is always convex and l.s.c. \cite{Rockafellar1998VariationalA}, it means that the $\text{prox}_{h^*}(\cdot)$ is often much easier to compute (potentially having a closed-form solution) than that of $h$.  For $h(\bz) = \sum_{i=1}^n \lambda ||z_i| - 1|$, we derive its conjugate function as illustrated in \Cref{fig:w_type_regularization}:
% \begin{align*}
% &h^*(\bx) = \sum_{i=1}^n h^*(x_i),\,\, \notag \\
% & \text{where } h^*(x_i) = \begin{cases}
%     |x_i|, & \text{if } x_i \in [-\lambda, \lambda],\\
%     +\infty, & \text{otherwise},
%     \end{cases}
% \end{align*} whose proximal operator  in the piecewise solution is
% \begin{align}
%     \label{Eq: the proximal mapping of conjugate function}
%     [\operatorname{prox}_{\tau h^*}(\by)]_i =
%     \begin{cases}
%         \lambda, & y_i > \lambda +\tau,\\
%         y_i -\tau, & \tau < y_i \leq \lambda +\tau, \\
%         0, & -\tau \leq y_i \leq\tau, \\
%         y_i +\tau, &  - \lambda -\tau \leq y_i < -\tau, \\
%         -\lambda, & y_i < -\lambda -\tau.
%     \end{cases}
% \end{align} This motivates us to perform a proximal step to update \(\bLmbd\).
% \begin{figure}[!t]
% \centering
% \includegraphics[width=\columnwidth]{papers_img/W_type_with_its_conjugate_func.pdf}
% \caption{W-type regularization function $h(z) = \lambda ||z| - 1|$ and its conjugate function $h^*(x)$.}
% \label{fig:w_type_regularization}
% \vspace{-1ex}
% \end{figure} We now outline \textbf{DualHash}, which employs an alternating update scheme for $\bx, \bB, \bm{\Lambda}$.
% we adopt a stochastic first-order method for the $\bx$-update since computing the full gradient \( \nabla_{\bx} F(\bx, \bB) \) is computationally expensive for typically the large number of samples $n$ in deep learning. Specifically, at each iteration $k$, we uniformly sample a mini-batch \(\calJ_k \subseteq \{1,\cdots,n\}\) with \(b_k = |\calJ_k| \leq n\) and compute:
% \begin{align}
% \label{Eq: the mini batch stochastic gradient estimator of F at x}
% \nabla_{\bx} F(\bx^k, \bB^k; \calJ_k)
% =\, \frac{1}{|\mathcal{J}_k|}\sum_{j \in \mathcal{J}_k} \nabla_{\bx} F_j(\bx^k, \bb_j^k).
% \end{align}
% Then we utilize this approximation within an optimization strategy chosen to enhance efficiency (e.g., incorporating momentum or variance reduction), as detailed in \Cref{Sec: Algorithm implementations}.
% Second, the auxiliary variable $\bB$ is updated via a simple gradient descent step \eqref{Eq: the update of B}; this suffices because the objective \eqref{Eq: simplified_lagrangian} with respect to $\bB$ is smooth with an inexpensive gradient,
% % which is slightly different from \cite{chambolleFirstOrderPrimalDualAlgorithm2011}:
%     \begin{align}
%    \nonumber \bB^{k+1} &= \bB^k - \tau \nabla_{\bB}\calL(\bx^{k+1}, \bB^k, \bLmbd^k)\\
%     &= \bB^k - \tau \left( \nabla_{\bB} F(\bx^{k+1}, \bB^k) + \bLmbd^k \right). \label{Eq: the update of B}
%     \end{align}
%     Finally, the dual variable $\bLmbd$ is updated via proximal gradient ascent on $h^*$, using an extrapolation step introduced in \cite{chambolleFirstOrderPrimalDualAlgorithm2011}. The update solves the following optimization problem:
% \begin{align}
%     \bLmbd^{k+1} &= \text{prox}_{(1/\tau h^*)} \left( \bLmbd^k + \left(2\bB^{k+1} - \bB^k  \right) \right) \notag \\
%     &=  \arg\min_{\bLmbd \in \RR^{n\times d} } \left\{ h^*(\bLmbd) - \langle \bLmbd, 2\bB^{k+1} - \bB^k \rangle \right. \notag \\
%     &\quad\quad\quad\quad\quad\quad\quad\quad\quad\quad\left. + \frac{\tau}{2} \| \bLmbd - \bLmbd^k\|^2 \right\}.
%     \label{Eq: the update of lambda}
% \end{align}
% \subsection{Implementations}
% \label{Sec: Algorithm implementations}
% The unified DualHash framework allows for flexible choices in the $\bx$ update step based on a mini-batch stochastic gradient approximation \eqref{Eq: the mini batch stochastic gradient estimator of F at x}. This flexibility enables us to incorporate advanced stochastic optimization techniques. We thus instantiate the $\bx$-update using two well-established approaches: \textbf{momentum} \cite{Nesterov1983AMF, Polyak1964} and \textbf{variance reduction} \cite{driggs2022biased}, which can improve empirical training efficiency and can also improve the theoretical convergence guarantees \cite{Liu2020AnIA,sutskever2013importance}.

% % \vspace{5pt}
% \paragraph{DualHash-StoM}
% %Numerically,
% The heavy ball method \cite{Polyak1964} and Nesterov acceleration \cite{Nesterov1983AMF} can be integrated with stochastic gradient descent (SGD) to yield accelerated variants like SGD with momentum, which significantly speeds up convergence.
% This approach has become a popular strategy for training neural networks \cite{fotopoulos2024review,Hertrich2020InertialSP}. We thus incorporate this technique and compute the stochastic gradient estimator (Line 3 of \Cref{alg:our_algorithm_sgdm}) as:
% \begin{align}
% \label{Eq: the stochastic approximate in SGDM}
%     \bG^k = \nabla_{\bx} F(\bz^k, \bB^k; \calJ_k),\, \text{where } \bz^k = \bx^k + \beta_k (\bx^k - \bx^{k-1}).
% \end{align} The update of \( \bx^k\) (Line 4 of Alg. \ref{alg:our_algorithm_sgdm}) is then:
% \begin{align}
% \label{Eq: the update of x in SGDM}
%     \bx^{k+1} = \by^k - \eta_k \bG^k,\,\text{where } \by^k = \bx^k + \alpha_k (\bx^k - \bx^{k-1}).
% \end{align} The overall algorithm, DualHash-StoM, is detailed in \Cref{alg:our_algorithm_sgdm}.
% \begin{remark}
% The momentum acceleration framework in \eqref{Eq: the update of x in SGDM} covers these two accelerations:
% Nesterov acceleration \cite{Nesterov1983AMF} corresponds to $\alpha_k = \beta_k \neq 0$; the heavy-ball method \cite{Polyak1964} is obtained with $\beta_k = 0$; vanilla SGD is recovered when $\alpha_k = \beta_k = 0$; In this paper, we set $\alpha_k, \beta_k \in (0,1)$ for all $k \geq 1$.
% \end{remark}

% % \vspace{3pt}
% \paragraph{DualHash-StoRM}
% Alternatively, we employ the STORM estimator \cite{cutkoskyMomentumBasedVarianceReduction2019} for the $\bx$-update. STORM utilizes a recursive momentum mechanism to reduce gradient variance, which is motivated by SGDM. Unlike some other variance reduction techniques (e.g., SVRG \cite{johnsonAcceleratingStochasticGradient2013}, SARAH \cite{SARAH2017}) that require full-gradient or large-batch gradient computations, STORM achieves variance reduction without such requirements. Given its practical effectiveness in deep learning \cite{cutkoskyMomentumBasedVarianceReduction2019,levy2021storm+}, we compute the stochastic gradient estimator (Line 3 of \Cref{alg:our_algorithm_storm}) as:
% % {\small
% % \begin{align}
% % \label{Eq: the stochastic approximate in STORM}
% % \bD^k =
% % \begin{cases}
% % \nabla_{\bx} F(\bx^1, \bB^1; \calJ_1), & \text{if } k = 1, \\
% % \nabla_{\bx} F(\bx^k, \bB^k; \calJ_k) + (1 - \rho_{k-1}) \left( \bD^{k-1} - \nabla_{\bx} F(\bx^{k-1}, \bB^{k-1};\calJ_{k-1})\right), & \text{if }k > 1.
% % \end{cases}
% % \end{align}}
% {\small
% \begin{align}
% \label{Eq: the stochastic approximate in STORM}
% \bD^k =
% \begin{cases}
% \nabla_{\bx} F(\bx^1, \bB^1; \mathcal{J}_1), &\text{if }  k = 1, \\
% (1 - \rho_{k-1}) \big( \bD^{k-1} - \nabla_{\bx} F(\bx^{k-1}, \bB^{k-1};\mathcal{J}_{k-1})\big)\\
% +\nabla_{\bx} F(\bx^k, \bB^k; \mathcal{J}_k), & \text{if } k > 1.
% \end{cases}
% \end{align}}
% The update of $\bx$ (Line 4 of \Cref{alg:our_algorithm_storm}) is then
% \begin{align}
% \label{Eq: the update of x in STORM}
%     \bx^{k+1} = \bx^k - \eta_k \bD^k.
% \end{align} The overall algorithm, DualHash-StoRM, is detailed in \Cref{alg:our_algorithm_storm}.

% \begin{figure}[!htbp]
%     \vspace{-5mm}
%     \begin{minipage}[t]{0.485\textwidth}
%         \begin{algorithm}[H]
%        \caption{DualHash-StoM}
%         \small   % 或使用 \footnotesize 更小
%         \setlength{\baselineskip}{0.8\baselineskip}  % 减少行间
%       \label{alg:our_algorithm_sgdm}
%       \begin{algorithmic}[1]
%         \STATE \textbf{Initialize:} $\bx^1 = \bz^1 = \by^1 \in \RR^{d_{\bx}} $, \( \bB^1 \in \RR^{n\times d},\, \bLmbd^1\in \RR^{n\times d}\), parameters $\{\eta_k\} $, $\tau$, $b_k$, $\{\alpha_k, \beta_k \in (0,1)\}$
%         \FOR{$k = 1, 2, \ldots$}
%           \STATE Compute \(\bG^k\) by \eqref{Eq: the stochastic approximate in SGDM}
%           \STATE Update $\bx^{k+1}$ by \eqref{Eq: the update of x in SGDM}
%           \STATE Update $\bB^{k+1}$ by \eqref{Eq: the update of B}
%           \STATE Update $\bm{\Lambda}^{k+1}$ by \eqref{Eq: the update of lambda}
%         \ENDFOR
%         \STATE \textbf{Return:} $\bx_R$ where \blue{$R$ is uniformly sampled from $\{1, \ldots, T\}$}
%       \end{algorithmic}
%     \end{algorithm}
%   \end{minipage}
%   \hspace{0.5mm}
%   \begin{minipage}[t]{0.485\textwidth}
%     \begin{algorithm}[H]
%         \caption{ DualHash-StoRM}
%       \small
%       \setlength{\baselineskip}{0.8\baselineskip}
%       \label{alg:our_algorithm_storm}
%       \begin{algorithmic}[1]
%         \STATE  \textbf{Initialize:} $\bx^1  \in \RR^{d_{\bx}}$, \( \bB^1 \in \RR^{n\times d},\, \bLmbd^1\in \RR^{n\times d}\), parameters $\{\eta_k\} $, $\tau $, $b_k$, $\{\rho_k\in (0,1)\}$
%         \FOR{$k = 1, 2, \ldots$}
%         \STATE Compute $\bD^k$ by \eqref{Eq: the stochastic approximate in STORM}
%         \STATE Update $\bx^{k+1}$ by \eqref{Eq: the update of x in STORM}.
%         \STATE Update $\bB^{k+1}$ by \eqref{Eq: the update of B}.
%         \STATE Update $\bm{\Lambda}^{k+1}$ by \eqref{Eq: the update of lambda}.
%         \ENDFOR
%         \STATE \textbf{Return:} $\bx_R$ where \blue{$R$ is uniformly sampled from $\{1, \ldots, T\}$}
%       \end{algorithmic}
%     \end{algorithm}
%   \end{minipage}
%   \vspace{-1ex}
% \end{figure}
% \begin{remark}
% The update steps for $\bB$ and $\bm{\Lambda}$ imply the following optimality conditions, which may result in explicit closed-solution in deep hashing. The $\bB$-update in \eqref{Eq: the update of B} yields:
% \begin{align}
%     \label{Eq: the optimality condition of Bk}
%     \bB^{k} = \bB^{k-1} - \tau \left( \nabla_{\bB}F(\bx^{k}, \bB^{k-1}) + \bm{\Lambda}^{k-1} \right).
% \end{align}
% For the $\bm{\Lambda}$-update in \eqref{Eq: the update of lambda}, the proximal operator definition gives:
% \[
%     \tau \left( \bm{\Lambda}^{k} + \tau^{-1}(2\bB^{k} - \bB^{k-1}) - \bm{\Lambda}^{k} \right) \in \partial h^{*}(\bm{\Lambda}^{k}),
% \]
% which is equivalent to the existence of $\mathcal{G}_{\bm{\Lambda}}^{k} \in \partial h^{*}(\bm{\Lambda}^{k})$ such that:
% \begin{align}
%     \label{Eq: the optimality condition of lbda}
%     \mathcal{G}_{\bm{\Lambda}}^{k} = (2\bB^{k} - \bB^{k-1}) + \tau(\bm{\Lambda}^{k-1} - \bm{\Lambda}^{k}).
% \end{align}
% If the sequence \(\{ (\bx^k, \bB^k, \bm{\Lambda}^k) \}\) converges to \((\bx^*, \bB^*, \bm{\Lambda}^*)\), then by the outer semicontinuity of $\partial h^{*}(\cdot)$, we have $\bB^* \in \partial h^*(\bm{\Lambda}^*)$.
% % which corresponds to the necessary condition \eqref{Eq: the sufficient and necessary condition for dual} for the second relation in optimality condition \eqref{Eq: the sufficient and necessary condition}.
% This establishes \eqref{Eq: the sufficient and necessary condition for dual}, which is precisely the necessary condition for the second relation in \eqref{Eq: the sufficient and necessary condition}.
% \end{remark}

\subsection{\textcolor{blue}{Min--Max Reformulation and KKT Conditions}}
\label{Sec: first-order optimality conditions}
Consider the equivalent constrained formulation of Problem~\eqref{Prop: block-wise regularized problem}:
\begin{align}
\min_{\bx \in \RR^{\blue{q}}, \bB \in \RR^{n\times d}, \bY \in \RR^{n\times d}} &\quad F(\bx, \bB) + h(\bY) \notag \\
\text{s.t. } &\quad \bY = \bB, \notag
\end{align}
whose dual problem is
\begin{align}
\label{Prop: the dual problem}
\max_{\bLmbd \in \RR^{n\times d}} \min_{\bx \in \RR^{\blue{q}}, \bB \in \RR^{n\times d}, \bY \in \RR^{n\times d}} &\quad F(\bx, \bB) + \langle \bLmbd, \bB - \bY \rangle + h(\bY). \tag{D}
\end{align}\blue{Recalling the definition of the conjugate function $h^*$} and observing that the minimization over $\bY$ in Problem~\eqref{Prop: the dual problem} can be carried out independently, we simplify the Lagrangian in Problem~\eqref{Prop: the dual problem} to
\begin{equation}\label{Eq: simplified_lagrangian}
\calL(\bx,\bB,\bLmbd):=\Phi(\bx,\bB,\bLmbd)-h^*(\bLmbd),
\qquad\text{where}\quad
\Phi(\bx,\bB,\bLmbd):=F(\bx,\bB)+\langle\bB,\bLmbd\rangle.
\end{equation}
Closely related to Problem~\eqref{Prop: the dual problem} is the nonconvex--concave min--max problem associated with $\calL$, which is equivalent to the convexified problem, i.e.,
\begin{equation}\label{Eq: minmax problem}
\min_{\bx,\bB}\max_{\bLmbd}\calL(\bx,\bB,\bLmbd)
=\min_{\bx,\bB}\ F(\bx,\bB)+h^{**}(\bB).
\end{equation}
where the equality follows from $\sup_{\bLmbd}\{\langle\bB,\bLmbd\rangle-h^*(\bLmbd)\}=h^{**}(\bB)$. Here, the function $\calL$ is concave in $\bLmbd$ but, through $\mathcal{D}$, nonconvex in $(\bx,\bB)$ jointly. Consequently, a saddle point of $\calL$ may fail to exist, and even checking whether a point is a local minimizer of the nonconvex problem \eqref{Prop: block-wise regularized problem} is NP-hard in general~\cite{Murty1987NP}. We therefore seek points satisfying the first-order necessary optimality condition for Problem~\eqref{Eq: minmax problem}. Following~\cite{chambolleFirstOrderPrimalDualAlgorithm2011,
clasonAccelerationGlobalConvergence2019,
botBohm2023alternating}, we call
$(\tilde{\bx},\tilde{\bB},\tilde{\bLmbd})$ a
\emph{KKT point} of $\calL$ if
$\bm{0}\in\mathcal{K}(\tilde{\bx},\tilde{\bB},\tilde{\bLmbd})$,
where
\begin{equation}\label{Eq: KKT mapping}
\mathcal{K}(\bx,\bB,\bLmbd):=
\begin{pmatrix}
\nabla_{\bx}F(\bx,\bB)\\
\nabla_{\bB}F(\bx,\bB)+\bLmbd\\
\bB-\partial h^*(\bLmbd)
\end{pmatrix}
=
\begin{pmatrix}
\nabla_{\bx}\calL(\bx,\bB,\bLmbd)\\
\nabla_{\bB}\calL(\bx,\bB,\bLmbd)\\
-\partial_{\bLmbd}(-\calL)(\bx,\bB,\bLmbd)
\end{pmatrix}.
\end{equation}
Equivalently,
\begin{equation}\label{Eq: KKT conditions}
\nabla_{\bx}F(\tilde{\bx},\tilde{\bB})=\bm{0},
\qquad
\nabla_{\bB}F(\tilde{\bx},\tilde{\bB})
+\tilde{\bLmbd}=\bm{0},
\qquad
\tilde{\bB}\in\partial h^*(\tilde{\bLmbd}).
\end{equation}Notice that the last relation in \eqref{Eq: KKT conditions}
is the first-order condition for maximizing $\calL$ with
respect to $\bLmbd$, namely
$\bm{0}\in\partial_{\bLmbd}(-\calL)$.
In contrast, for standard Lagrangian-based methods such as
ADMM and ALM~\cite{wangGlobalConvergenceADMM2018,Bolte2018},
the multiplier enters affinely, so
$\bm{0}\in\partial\calL$ yields the usual KKT system,
with the multiplier component enforcing primal feasibility.
In our formulation, however, the nonsmooth concave term
$-h^*(\bLmbd)$ prevents identifying
$\bm{0}\in\partial\calL$ with
$\bm{0}\in\mathcal{K}$. We therefore use
$\bm{0}\in\mathcal{K}$. In the following, we denote by $\operatorname{crit}\calL:=\{(\bx, \bB, \bLmbd) \in \RR^{q} \times \RR^{n\times d} \times \RR^{n\times d}: \bm{0} \in \mathcal K(\bx, \bB, \bLmbd)\}$ the set of KKT points of $\calL$.

In fact, every local saddle point of $\calL$ is a KKT point by Fermat's rule~\cite[Theorem~10.1]{Rockafellar1998VariationalA}. Moreover, for every $(\tilde{\bx},\tilde{\bB},\tilde{\bLmbd})$ with $\tilde{\bLmbd}\in\dom h^*$, KKT points of $\calL$ correspond one-to-one to stationary points of the convexified problem~\eqref{Eq: minmax problem} through $\tilde{\bLmbd}=-\nabla_{\bB}F(\tilde{\bx},\tilde{\bB})$, i.e.,
\begin{equation}\label{Eq: saddle-KKT-convexified}
\begin{aligned}
\bm{0}\in\mathcal{K}(\tilde{\bx},\tilde{\bB},\tilde{\bLmbd})
\quad\Longleftrightarrow\quad
\tilde{\bLmbd}=-\nabla_{\bB}F(\tilde{\bx},\tilde{\bB})\ \text{ and }\ \bm{0}\in\nabla F(\tilde{\bx},\tilde{\bB})+\{\bm{0}\}\times\partial h^{**}(\tilde{\bB}).
\end{aligned}
\end{equation}
where the equivalence follows from the subgradient inversion and sum rules~\cite[Theorems~23.5 and~23.8]{rockafellar1970convex}. This observation motivates the following characterization of when a KKT point of the min--max formulation~\eqref{Eq: minmax problem} is also stationary for Problem~\eqref{Prop: block-wise regularized problem}.

\subsection{Exact Stationarity Relations for Problem~\eqref{Prop: block-wise regularized problem}}
\label{Sec: original optimality relations}
We first prove that any primal--dual optimal solution $(\bx^*,\bB^*,\bLmbd^*)$ with zero duality gap is a KKT point of $\calL$. The proof is provided in \Cref{app:zero-gap-proof}.
\begin{proposition}\label{prop:zerogap}
Suppose that $(\bx^*,\bB^*)$ and $\bLmbd^*$ are optimal solutions of Problems~\eqref{Prop: block-wise regularized problem} and~\eqref{Prop: the dual problem}, respectively, with zero duality gap. Then $(\bx^*,\bB^*,\bLmbd^*)\in\operatorname{crit}\calL$.
\end{proposition}
Recall that strong duality holds for Problems~\eqref{Prop: block-wise regularized problem} and~\eqref{Prop: the dual problem} if their optimal values coincide (zero duality gap) and both are attained. Here, since $J$ is coercive and using $\phi^*(y)\ge|y|$ for a W-type regularizer (Lemma~\ref{lem:wtype-conj}), both problems admit optimal solutions, so zero duality gap and strong duality are equivalent and hold or fail together. Let $p^*$, $p^{**}$, and $d^*$ denote the optimal values of Problem~\eqref{Prop: block-wise regularized problem}, of the right-hand side of~\eqref{Eq: minmax problem}, and of Problem~\eqref{Prop: the dual problem}, respectively. The duality gap splits as
\[
p^*-d^*=(p^*-p^{**})+(p^{**}-d^*),
\]
where the first term is the loss caused by replacing $h$ with $h^{**}$, and the second term stems from the nonconvexity of $\calL$ in $(\bx,\bB)$. The following example shows that, even when $F$ is convex so that $p^{**}=d^*$, the first term can be positive, which motivates relating KKT points of $\calL$ directly to Problem~\eqref{Prop: block-wise regularized problem}
\begin{example}\label{ex:gap}
Consider Problem~\eqref{Prop: block-wise regularized problem} with $n=d=1$ and without the $\bx$-block, and let $F(b)=\tfrac12 b^2$ and $\lambda=1$. The primal value is
$p^*=\min_{b\in\RR}\bigl\{\tfrac12 b^2+\bigl||b|-1\bigr|\bigr\}=\tfrac12$, attained at $b^*=\pm1$.
Since $\min_b\{\tfrac12 b^2+\Lambda b\}=-\tfrac12\Lambda^2$ and $h^*(\Lambda)=|\Lambda|$ on $[-1,1]$ ($+\infty$ elsewhere), the dual value is
\[
d^*=\sup_{|\Lambda|\le1}\Bigl\{-\tfrac12\Lambda^2-|\Lambda|\Bigr\}=0,
\]
attained at $\Lambda^*=0$. Hence $p^*-d^*=\tfrac12>0$, and no zero-gap primal--dual pair exists.
\end{example}

A zero-gap pair therefore need not exist, and we relate KKT points of $\calL$ directly to Problem~\eqref{Prop: block-wise regularized problem}, which holds for every W-type regularizer. The proof is provided in \Cref{app:transfer-proof}.
\begin{proposition}\label{prop:transfer}
Let \(h\) be a W-type regularizer, and let $(\tilde{\bx},\tilde{\bB},\tilde{\bLmbd})$ be a KKT point of $\calL$. The following statements are equivalent:
\begin{enumerate}[label=\textup{(\alph*)}]
\item $(\tilde{\bx},\tilde{\bB})$ is a stationary point of Problem~\eqref{Prop: block-wise regularized problem}, i.e.,
$\bm{0}\in\nabla F(\tilde{\bx},\tilde{\bB})+\{\bm{0}\}\times\partial h(\tilde{\bB})$;
\item $h(\tilde{\bB})+h^*(\tilde{\bLmbd})=\langle\tilde{\bLmbd},\tilde{\bB}\rangle$;
\item $|\tilde{b}_{ij}|\ge1$ for all $i,j$.
\end{enumerate}
\end{proposition}

\subsection{From KKT Residuals to Problem~\eqref{Eq: the primal deep hashing problem}}
\label{Sec: original objective guarantees}
Proposition~\ref{prop:transfer} characterizes when a KKT point
is stationary for the penalized problem.
From an optimization perspective, however, for the original
objective $J_0$, we must further account for
the constraint violation and the convexification error,
namely,
\[
\|\bB-\mathcal{D}(\bx)\|,
\qquad
h(\mathcal{D}(\bx))-h^{**}(\mathcal{D}(\bx)).
\]
In this subsection, Fenchel conjugacy allows us to control the effect of the
constraint violation through the KKT residual
$r_{\mathcal{K}}:=\operatorname{dist}\bigl(\bm{0},\mathcal{K}(\bx,\bB,\bLmbd)\bigr)$,
while the structure of the W-type regularizer bounds the convexification error by
$h(\mathcal{D}(\bx))$, which equals $\lambda Q_1(\bx)$ for the $\ell_1$ penalty, with $Q_1$ defined in~\eqref{eq:quantization-error}. Consequently, we establish
that the stationarity of $J_0$ at a point $\widehat{\bx}$ near $\bx$ is bounded by an
optimization error and a structural error:
\[
\operatorname{dist}\bigl(\bm{0},\partial J_0(\widehat{\bx})\bigr)
\le
\underbrace{\calO\bigl(r_{\mathcal{K}}\bigr)}_{\text{optimization error}}
+
\underbrace{\calO\bigl(\sqrt{h(\mathcal{D}(\bx))}\bigr)}_{\text{structural error}}.
\]

Before proceeding, due to the nonsmooth, nonconvex nature of the objective $J_0$, we measure its stationarity using a nearby-point criterion related to near-approximate stationarity~\cite[Definition~2.5]{tian2022stationarity} as follows.
\begin{definition}[Near-stationarity]\label{def:near-stationarity}
Given $\delta,\varepsilon\ge0$, a point $\bx$ is called \emph{$(\delta,\varepsilon)$-near-stationary} for Problem~\eqref{Eq: the primal deep hashing problem} if there exists $\widehat{\bx}$ such that
\begin{equation}\label{eq:near-stationarity}
\|\widehat{\bx}-\bx\|\le\delta,
\qquad
\operatorname{dist}\bigl(\bm{0},\partial J_0(\widehat{\bx})\bigr)\le\varepsilon,
\end{equation}
where $\partial$ denotes the limiting subdifferential defined in~\eqref{Eq: limiting subdifferential}.
\end{definition}
For $\bLmbd\in\operatorname{dom}h^*$, let us define the potential function used in this subsection,
\[
s_{\bLmbd}(\by):=f(\by)+\langle\bLmbd,\mathcal{D}(\by)\rangle-h^*(\bLmbd).
\]
By Assumption~\ref{Ass: basic regularity}(ii), $s_{\bLmbd}$ is
$L_{s_{\bLmbd}}:=(L_f+\|\bLmbd\|L_{\nabla\mathcal{D}})$-smooth.
If $\operatorname{dom}h^*$ is bounded, as for $p=1$, log-cosh, and Huber in
\Cref{tab:wtype}, then $\operatorname{dom}\phi^*=[-\lambda,\lambda]$ gives
$\|\bLmbd\|\le\lambda\sqrt{nd}$, and hence
\[
L_{s_{\bLmbd}}\le L_0:=L_f+\lambda\sqrt{nd}\,L_{\nabla\mathcal{D}},
\]
a constant independent of $\bLmbd$.
For notational simplicity, we set
$m:=\gamma/n$ and $\kappa:=(1+m^2)/(4m)$. The following lemma shows that, at $\bx$, the
first-order error of $s_{\bLmbd}$ is governed by the KKT residual, while its gap to $J_0$
is governed by the network outputs and is never smaller than the convexification error;
the structural error is therefore intrinsic to passing from the convexified model \eqref{Eq: minmax problem} to the original problem \eqref{Prop: block-wise regularized problem}.

\begin{lemma}\label{lem:conjugate-minorant}
\red{Under Assumption~\ref{Ass: basic regularity}(ii),} let $h$ be a W-type
regularizer defined in \Cref{def:wtype}, and let $(\bx,\bB,\bLmbd)$ satisfy
$r_{\mathcal{K}}<\infty$. \red{Then $s_{\bLmbd}\le J_0$ on $\RR^q$, and the following
inequalities hold at the given point $\bx$:}
\begin{align}
\|\nabla s_{\bLmbd}(\bx)\|
&\le\sqrt{1+G_{\mathcal{D}}^2}\;r_{\mathcal{K}},
\label{eq:minorant-gradient}\\
0\le h(\mathcal{D}(\bx))-h^{**}(\mathcal{D}(\bx))
\le J_0(\bx)-s_{\bLmbd}(\bx)
&\le h(\mathcal{D}(\bx))+\kappa\,r_{\mathcal{K}}^2.
\label{eq:minorant-gap}
\end{align}
In particular, if $\mathcal{D}(\bx)\in[-1,1]^{n\times d}$, it holds that
$$0 \le h(\mathcal{D}(\bx))\le J_0(\bx)-s_{\bLmbd}(\bx)\le h(\mathcal{D}(\bx))+\kappa\,r_{\mathcal{K}}^2.$$
\end{lemma}
We are now ready to bound the near-stationarity of $J_0$ explicitly. The proof is given in \Cref{app:original-certificate-proof}.
\begin{theorem}\label{prop:original-certificate}
Under the conditions of \Cref{lem:conjugate-minorant},
let $L>0$ satisfy
\red{$L\ge L_s(\|\bLmbd\|)$.} Then there exists $\widehat{\bx}\in\RR^q$ such that
\begin{align*}
\|\widehat{\bx}-\bx\|
&\le\sqrt{\frac{2}{L}}
\Bigl(\sqrt{\kappa}\,r_{\mathcal{K}}
+\sqrt{h(\mathcal{D}(\bx))}\Bigr),\\
\operatorname{dist}\bigl(\bm{0},\partial J_0(\widehat{\bx})\bigr)
&\le
\underbrace{
\Bigl(\sqrt{1+G_{\mathcal{D}}^2}
+2\sqrt{2\kappa L}\Bigr)r_{\mathcal{K}}
}_{\text{optimization error}}
+\underbrace{
2\sqrt{2L\,h(\mathcal{D}(\bx))}
}_{\text{structural error}}.
\end{align*}
\red{Moreover, if $h$ has a compact dual domain, one may take
$L=L_0:=L_f+\lambda\sqrt{nd}\,L_{\nabla\mathcal{D}}$,
so that the bounds hold uniformly over $\bLmbd$.}
\end{theorem}
The theorem connects the near-stationarity of the original objective $J_0$ to two quantities: the KKT residual, which motivates our oracle-complexity analysis, and the regularization term $h(\mathcal{D}(\bx))$, which is determined by the network outputs. For the $\ell_1$ W-type regularizer, the latter equals $\lambda Q_1(\bx)$; thus, the quantization error contributes directly to the near-stationarity bound and motivates our use of the $\ell_1$ W-type regularizer in the experiments.

\begin{remark}[Role of the conjugate minorant]\label{rem:conjugate-minorant-role}
\blue{Since $h\ge0$, the loss $f$ is also a smooth minorant of $J_0$, but its gradient $\nabla f(\bx)=\nabla s_{\bLmbd}(\bx)-\nabla\mathcal{D}(\bx)^*\bLmbd$ retains the multiplier in the first-order term. The minorant $s_{\bLmbd}$ absorbs $\langle\bLmbd,\mathcal{D}(\cdot)\rangle$, so the multiplier enters \Cref{prop:original-certificate} only through the curvature constant $L_0$. Consequently, the structural error does not depend on the penalty parameter $\gamma$, which enters only the optimization error through $\kappa$; the certificate holds for every fixed $\gamma>0$, without the limit $\gamma\to\infty$ required by classical quadratic penalty methods~\cite[Sec.~17.1]{nocedalWright2006}.}
\end{remark}

